\documentclass[final]{beamer}
\usepackage[T1]{fontenc}
\usepackage[utf8]{inputenc}
\usepackage[english]{babel}
\usepackage[size=a0,scale=1.1]{beamerposter}
\usepackage{amsmath,amssymb}
\usepackage{booktabs}
\usepackage{tabularx}
\usepackage{qrcode}
\usepackage{tikz}
\usepackage{graphicx}
\usepackage{multicol}
\usepackage{ragged2e}
\usepackage{microtype}
\usepackage{setspace}
\usetikzlibrary{arrows.meta,positioning,shapes.geometric}
\microtypesetup{expansion=false}
\newcolumntype{C}{>{\centering\arraybackslash}X}

% ---------- Colors ----------
\definecolor{PosterPrimary}{RGB}{24,40,60}
\definecolor{PosterSecondary}{RGB}{40,84,104}
\definecolor{PosterAccent}{RGB}{78,164,184}
\definecolor{PosterSurface}{RGB}{249,251,253}
\definecolor{PosterMuted}{RGB}{234,242,246}
\definecolor{PosterText}{RGB}{32,36,45}
\definecolor{PosterGreen}{RGB}{88,196,65}

% ---------- Beamer style ----------
\setbeamercolor{normal text}{fg=PosterText,bg=white}
\setbeamercolor{headline}{fg=white,bg=PosterPrimary}
\setbeamercolor{block title}{fg=white,bg=PosterSecondary}
\setbeamercolor{block body}{fg=PosterText,bg=PosterSurface}
\setbeamercolor{block alerted title}{fg=white,bg=PosterAccent}
\setbeamercolor{block alerted body}{fg=PosterText,bg=PosterMuted}
\setbeamercolor{item}{fg=PosterSecondary}
\setbeamercolor{itemize item}{fg=PosterSecondary}
\setbeamercolor{itemize subitem}{fg=PosterSecondary}

\setbeamerfont{headline title}{size=\VeryHuge,series=\bfseries}
\setbeamerfont{headline author}{size=\LARGE}
\setbeamerfont{headline institute}{size=\Large}
\setbeamerfont{block title}{size=\large,series=\bfseries}
\setbeamerfont{block body}{size=\normalsize,series=\mdseries}
\setbeamerfont{itemize/enumerate body}{size=\normalsize}
\setbeamertemplate{itemize item}{\raise0.15em\hbox{\scriptsize$\blacktriangleright$}}
\setbeamertemplate{itemize subitem}{\raise0.2em\hbox{\tiny$\bullet$}}
\setbeamertemplate{navigation symbols}{}
\setbeamertemplate{blocks}[rounded][shadow=false]
\setlength{\leftmargini}{1.5em}
\setlength{\leftmarginii}{1.1em}

\setlength{\belowcaptionskip}{1.0ex}
\setlength{\abovecaptionskip}{1.0ex}
\setlength{\parskip}{0.55em}
\setlength{\columnsep}{2.2cm}
\setbeamersize{text margin left=1.4cm,text margin right=1.4cm}
\setstretch{1.14}
\addtobeamertemplate{block begin}{}{\vspace{0.5em}}
\addtobeamertemplate{block end}{\vspace{0.35em}}{}

\setbeamertemplate{headline}{
  \leavevmode
  \begin{beamercolorbox}[wd=\paperwidth,ht=9.4cm,dp=1.4cm,leftskip=1.8cm,rightskip=1.8cm]{headline}
    \centering
    {\usebeamerfont{headline title}\inserttitle\par}
    \vspace{0.35em}
    {\usebeamerfont{headline author}\insertauthor\par}
    \vspace{0.2em}
    {\usebeamerfont{headline institute}\insertinstitute}
  \end{beamercolorbox}
  \begin{beamercolorbox}[wd=\paperwidth,ht=0.42cm,dp=0cm]{block alerted title}
  \end{beamercolorbox}
}

\setbeamertemplate{footline}{
  \leavevmode
  \begin{beamercolorbox}[wd=\paperwidth,ht=1.2cm,dp=0.4cm,leftskip=2.4cm,rightskip=2.4cm]{normal text}
    \color{gray}\small \insertshorttitle\hfill\insertshortauthor
  \end{beamercolorbox}
}

% ---------- Metadata ----------
\title{Domain-Specific Tokenization for Biomedical NLP}
\author{Anna Belyakova}
\institute{Innopolis University}
\date{\today}

\begin{document}
\begin{frame}[t]

\begin{columns}[t,totalwidth=\textwidth]

% ========== Column 1 ==========
\begin{column}{0.32\textwidth}

\begin{block}{Project Access}
\begin{flushleft}
\setlength{\tabcolsep}{0.4em}
\renewcommand{\arraystretch}{1.0}
\begin{tabular}{@{}p{3.2cm}p{8.8cm}@{}}
\fcolorbox{PosterSecondary}{white}{\qrcode[height=2.9cm]{https://github.com/belyakova-anna/Domain-specific-Tokenization}}
&
\begin{minipage}[t]{8.6cm}
\vspace*{-2em}
\setlength{\parskip}{0pt}
\setstretch{1.0}
{\normalsize\textbf{Repository:}}\\[0.1em]
{\small notebooks + tokenizers + poster code}\\[0.3em]
{\normalsize\textbf{Study topic:}}\\[0.1em]
{\small domain-specific tokenization}
\end{minipage}
\end{tabular}
\end{flushleft}
\end{block}

\begin{block}{Motivation}
\justifying
\begin{itemize}
    \item Biomedical text contains long compounds, abbreviations, and rare terms that are hard for generic tokenizers.
    \item Over-fragmentation increases sequence length, weakens entity boundaries, and may harm minority classes.
    \item Goal: train compact domain BPE tokenizers and test whether improved segmentation translates to downstream gains.
\end{itemize}
\end{block}

\begin{block}{Dataset Signals (PubMedQA)}
\justifying
\begin{center}
\begin{tikzpicture}[x=1cm,y=1cm]
\node[
  draw=PosterSecondary,
  rounded corners=8pt,
  fill=white,
  minimum width=14.8cm,
  minimum height=7.5cm,
  anchor=north west
] (card) at (0,0) {};
\node[
  anchor=north west,
  text width=11.4cm,
  align=left,
  font=\footnotesize
] at (0.55,-0.55) {%
{\bfseries\large Corpus Profile}\par\vspace{0.35em}
\textbf{Labeled:} 1,000 \hspace{1.0em} \textbf{Unlabeled:} 61,249\par\vspace{0.28em}
\textbf{Question length:} mean 12.9 words\par\vspace{0.28em}
\textbf{Context length:} mean 200.2 words\par\vspace{0.28em}
\mbox{\textbf{Class balance:} yes 55.2\%, no 33.8\%, maybe 11.0\%}
};
\end{tikzpicture}
\end{center}
\vspace{0.2em}
\begin{itemize}
    \item Contexts dominate tokenization cost and drive model efficiency.
    \item Minority class (\textit{maybe}) motivates macro-F1 over accuracy-only evaluation.
\end{itemize}
\end{block}

\begin{block}{Experimental Setup}
\justifying
\begin{flushleft}
\scriptsize\bfseries Baselines
\end{flushleft}
\begin{flushleft}
\begin{tabular}{@{}lll@{}}
\tikz[baseline=(n.base)]\node[draw=PosterSecondary,rounded corners=7pt,fill=white,minimum width=3.25cm,minimum height=1.1cm,inner sep=2pt,font=\small\bfseries] (n) {BERT-base}; &
\hspace{0.4cm}\tikz[baseline=(n.base)]\node[draw=PosterSecondary,rounded corners=7pt,fill=white,minimum width=3.25cm,minimum height=1.1cm,inner sep=2pt,font=\small\bfseries] (n) {BioBERT}; &
\hspace{0.4cm}\tikz[baseline=(n.base)]\node[draw=PosterSecondary,rounded corners=7pt,fill=white,minimum width=3.25cm,minimum height=1.1cm,inner sep=2pt,font=\small\bfseries] (n) {GPT-2}; \\
\end{tabular}
\end{flushleft}
\vspace{0.35em}
\begin{flushleft}
\scriptsize\bfseries Custom
\end{flushleft}
\begin{flushleft}
\begin{tabular}{@{}ll@{}}
\tikz[baseline=(n.base)]\node[draw=PosterAccent,rounded corners=7pt,fill=PosterMuted,minimum width=3.75cm,minimum height=1.1cm,inner sep=2pt,font=\small\bfseries] (n) {BioBPE-30k}; &
\hspace{0.5cm}\tikz[baseline=(n.base)]\node[draw=PosterAccent,rounded corners=7pt,fill=PosterMuted,minimum width=3.75cm,minimum height=1.1cm,inner sep=2pt,font=\small\bfseries] (n) {BioBPE-50k}; \\
\end{tabular}
\end{flushleft}
\vspace{0.2em}
\small
\textbf{Protocol:} 300 samples for intrinsic analysis; PubMedQA 70/15/15 for classification; JNLPBA 4k/1k for NER metrics.
\end{block}

\begin{alertblock}{Hypothesis}
\justifying
Domain-trained BPE will reduce fragmentation in biomedical text and improve task-level robustness, especially for entity-heavy inputs.
\end{alertblock}

\end{column}

% ========== Column 2 ==========
\begin{column}{0.32\textwidth}

\begin{block}{Tokenizer Adaptation Pipeline}
\justifying
\begin{center}
\begin{tikzpicture}[
    node distance=1.15cm,
    box/.style={draw=PosterSecondary,fill=white,rounded corners=8pt,line width=1.4pt,minimum width=8.8cm,minimum height=1.15cm,align=center,font=\bfseries},
    >=Latex
]
\node[box] (a) {1) Collect PubMedQA corpora};
\node[box,below=of a] (b) {2) Train biomedical BPE vocabularies (30k/50k)};
\node[box,below=of b] (c) {3) Evaluate intrinsic metrics (fertility, compactness, UNK)};
\node[box,below=of c] (d) {4) Validate on classification and NER-oriented metrics};
\node[box,below=of d] (e) {5) Rank tokenizers by efficiency + quality trade-off};
\draw[->,line width=1.8pt,PosterAccent] (a) -- (b);
\draw[->,line width=1.8pt,PosterAccent] (b) -- (c);
\draw[->,line width=1.8pt,PosterAccent] (c) -- (d);
\draw[->,line width=1.8pt,PosterAccent] (d) -- (e);
\end{tikzpicture}
\end{center}
\vspace{0.2em}
\justifying
\textbf{Qualitative term behavior:} baseline tokenizers split terms like \textit{gastroesophageal} into 5+ pieces, while custom BPE often keeps the whole term in one token.
\end{block}

\begin{block}{Tokenizer Adaptation Methods}
\justifying
\begin{beamercolorbox}[wd=\linewidth,rounded=true,sep=0.65em]{block body}
\footnotesize\textbf{FVT (Fast Vocabulary Transfer)}
\[
\footnotesize
E_{\text{new}}(t_i)=\frac{1}{\lvert T_d(t_i)\rvert}\sum_{t_j\in T_d(t_i)}E_{\text{old}}(t_j)
\]
\small Initializes each new token embedding as the average of embeddings of its source subpieces from the old vocabulary.
\end{beamercolorbox}
\vspace{0.25em}
\begin{beamercolorbox}[wd=\linewidth,rounded=true,sep=0.65em]{block body}
\footnotesize\textbf{FOCUS}
\[
\footnotesize
E_{\text{new}}(t_i)=\frac{1}{\lvert V_t\rvert}\sum_{t_j\in V_t}w_t\,E_{\text{old}}(t_j)
\]
\small Uses weighted averaging with similarity-aware coefficients to better preserve semantic neighborhoods.
\end{beamercolorbox}
\vspace{0.25em}
\begin{beamercolorbox}[wd=\linewidth,rounded=true,sep=0.65em]{block body}
\footnotesize\textbf{ZeTT (Zero-shot Tokenizer Transfer)}
\[
\footnotesize
\mathcal{L}_{\phi}=\mathcal{L}_{\text{LM}}+\alpha\cdot\mathcal{L}_{\text{aux}}
\]
\small Learns tokenizer-adaptation parameters by combining language modeling loss with tokenizer-specific auxiliary constraints.
\end{beamercolorbox}
\end{block}

\begin{block}{Evaluation Signals}
\justifying
\begin{itemize}
    \item \textbf{Segmentation compactness:} avg tokens per sample, fertility, token length.
    \item \textbf{Downstream utility:} test accuracy and macro-F1 on yes/no/maybe.
    \item \textbf{NER compatibility:} subtoken-per-entity ratio, single-piece entity rate, fragmentation rate.
\end{itemize}
\vspace{0.2em}
\textbf{Composite NER tokenization score}
\[
\text{Score}=0.50(1-\text{Frag})+0.30(\text{SinglePiece})+0.20(1-\text{UNK})
\]
\end{block}

\end{column}

% ========== Column 3 ==========
\begin{column}{0.32\textwidth}

\begin{block}{Experimental Results}
\justifying
\textbf{Stage 4: Intrinsic tokenization metrics (best to worst by fertility)}
\vspace{0.25em}
\begin{center}
\scriptsize
\begin{tabularx}{0.96\linewidth}{lCCC}
\toprule
Tokenizer & Avg Tok/Sample & Fertility & Avg Tok Len \\
\midrule
BioBPE-50k & 273.26 & \textbf{1.301} & 5.12 \\
BioBPE-30k & 275.11 & 1.310 & 5.08 \\
GPT-2 & 313.38 & 1.492 & 5.01 \\
BERT-base & 323.89 & 1.542 & 4.29 \\
BioBERT & 339.49 & 1.616 & 4.05 \\
\bottomrule
\end{tabularx}
\end{center}
\vspace{0.15em}
\textbf{Relative gain:} best custom model reduces fertility by \(\sim\)12.8\% vs best baseline (GPT-2).
\vspace{0.55em}

\textbf{Stage 5: Downstream classification (test set)}
\begin{center}
\scriptsize
\begin{tabularx}{0.96\linewidth}{lCC}
\toprule
Tokenizer & Accuracy & Macro-F1 \\
\midrule
BioBPE-50k & \textbf{0.5667} & \textbf{0.3718} \\
BioBPE-30k & \textbf{0.5667} & 0.3710 \\
GPT-2 & \textbf{0.5667} & 0.3643 \\
BERT-base & 0.5600 & 0.3582 \\
BioBERT & 0.5533 & 0.3521 \\
\bottomrule
\end{tabularx}
\end{center}
\vspace{0.2em}
\textbf{Observation:} custom BPE improves macro-F1 while keeping accuracy stable.
\vspace{0.5em}

\textbf{Stage 6: NER-oriented ranking}
\begin{center}
\begin{tikzpicture}
\node[anchor=west] at (-1.0,2.45) {\small BioBPE-50k};
\fill[PosterSecondary!80] (2.9,2.2) rectangle (8.0,2.6);
\node[anchor=west] at (-1.0,1.65) {\small BioBPE-30k};
\fill[PosterSecondary!65] (2.9,1.4) rectangle (7.55,1.8);
\node[anchor=west] at (-1.0,0.85) {\small BERT-base};
\fill[PosterAccent!75] (2.9,0.6) rectangle (6.45,1.0);
\node[anchor=west] at (-1.0,0.05) {\small GPT-2};
\fill[PosterAccent!60] (2.9,-0.2) rectangle (6.2,0.2);
\node[anchor=west] at (-1.0,-0.75) {\small BioBERT};
\fill[PosterAccent!45] (2.9,-1.0) rectangle (6.0,-0.6);
\draw[PosterText!35,line width=0.8pt] (2.9,-1.25) -- (8.9,-1.25);
\node[anchor=west,font=\scriptsize] at (2.9,-1.6) {NER tokenization score};
\node[anchor=west,font=\scriptsize,fill=white,inner sep=1.2pt] at (8.15,2.4) {0.3049};
\node[anchor=west,font=\scriptsize,fill=white,inner sep=1.2pt] at (7.70,1.6) {0.2920};
\node[anchor=west,font=\scriptsize,fill=white,inner sep=1.2pt] at (6.60,0.8) {0.2272};
\node[anchor=west,font=\scriptsize,fill=white,inner sep=1.2pt] at (6.35,0.0) {0.2167};
\node[anchor=west,font=\scriptsize,fill=white,inner sep=1.2pt] at (6.15,-0.8) {0.2087};
\end{tikzpicture}
\end{center}
\end{block}

\begin{alertblock}{Key Takeaways}
\begin{itemize}
  \item Best tokenizer overall: \textbf{BioBPE-50k}.
  \item Intrinsic compactness and extrinsic quality improve consistently.
  \item Gains are strongest in entity-sensitive biomedical text.
\end{itemize}
\end{alertblock}

\begin{block}{Conclusions}
\justifying
\begin{itemize}
    \item Domain-adapted BPE produces consistently more compact segmentation on biomedical text.
    \item Custom tokenizers deliver the best macro-F1 and strongest NER-oriented scores.
    \item \(50k\) vocabulary is the most stable option across intrinsic and extrinsic evaluations.
    \item Practical recommendation: use custom biomedical BPE as a default tokenizer for PubMed-style pipelines.
\end{itemize}
\textbf{Limitations:} fixed lightweight classifier and single-domain corpus; future work includes transformer fine-tuning with embedding transfer (FVT/FOCUS/ZeTT).
\end{block}

\begin{block}{Reproducibility}
\justifying
\begin{itemize}
    \item Notebooks: 6 staged experiments from EDA to NER tokenization evaluation.
    \item Deterministic setup: fixed seeds and consistent data protocol.
    \item All metrics reported on the same benchmark splits for fair comparison.
\end{itemize}
\end{block}

\end{column}

\end{columns}
\end{frame}
\end{document}
